\section{Related Work}
\label{sec:related}

\paragraph{Specification gaming and reward hacking.}
Agents that optimize an imperfect reward often find behaviors the designer did
not intend \citep{amodei2016concrete,krakovna2020specification}.
\citet{skalse2022rewardhacking} formalize reward hacking as a proxy whose
optimization decreases the true return. \citet{pan2022effects} show that more
capable agents exploit misspecified rewards more, sometimes through abrupt phase
transitions, and \citet{gao2023scaling} measure how a true objective first
improves and then degrades as a learned reward model is optimized. Our dose-
response and training-dynamics experiments (Sections~\ref{sec:dynamics} and
\ref{sec:dose}) ask the same questions in settings small enough that the true
objective is known exactly and the misspecification can be switched off.

\paragraph{Reward tampering.}
\citet{everitt2021reward} distinguish an agent influencing the process that
computes its reward from merely exploiting a flawed reward function.
\citet{denison2024sycophancy} find that language models trained on a curriculum
of gameable environments occasionally generalize to editing their own reward
function. \ST{} is a minimal instance of the first kind: the reward is computed
from a delivery sensor, and the agent can make that sensor report deliveries
that did not happen.

\paragraph{Shaping and goal misgeneralization.}
Shaping rewards preserve optimal policies only if they are potential-based
\citep{ng1999policy}; \PF{} uses a non-potential bonus precisely so that it
creates a farmable cycle. \citet{langosco2022goal} show that agents can keep
their capabilities under distribution shift while pursuing the wrong goal. Our
OOD splits are simpler layout shifts (more walls, a more distant console, more
gems); they test whether a learned exploit persists, not whether goals
misgeneralize.

\paragraph{Monitoring for misbehavior.}
\citet{pan2022effects} propose detecting aberrant policies as an anomaly-
detection task. \citet{baker2025monitoring} show that monitors reading a
reasoning model's chain of thought catch reward hacking in agentic coding, and
that optimizing against such a monitor can teach the model to hide its intent.
\citet{macdiarmid2025natural} report that reward hacking learned in production
RL environments generalizes to broader misalignment, which raises the stakes of
detecting it early. Our monitors see only trajectories, but the evaluation
questions are the same: what may the monitor observe, does it generalize to new
conditions, and is it better than a trivial baseline?

\paragraph{Small testbeds and evaluation methodology.}
AI Safety Gridworlds \citep{leike2017ai} and bsuite \citep{osband2020bsuite}
argue for small, diagnostic environments; Procgen \citep{cobbe2020leveraging}
argues for procedurally generated layouts so that evaluation is not
memorization. \citet{agarwal2021precipice} show how few-seed RL comparisons
mislead and recommend interval estimates over runs. We follow all three
recommendations and add one of our own: every number in the paper is generated
from saved results and checked in continuous integration
(Section~\ref{sec:pitfalls}).
